\documentclass[11pt]{article}
% math symbols
\usepackage{amssymb,amsmath}
% for different compilers
\usepackage{ifpdf}
% geometry of page
\usepackage[margin=2.1cm]{geometry}
% float pictures
\usepackage{wrapfig}

% if pdflatex, then
\ifpdf
 \usepackage[english,russian]{babel}
 \usepackage[utf8]{inputenc}
 \usepackage[unicode]{hyperref}
 \usepackage[pdftex]{graphicx}
 \usepackage{cmlgc}
% if xelatex, then
\else
% math fonts
 \usepackage{fouriernc}
% xelatex specific packages
 \usepackage[xetex]{hyperref}
 \usepackage{xunicode}	% some extra unicode support
 \usepackage{xltxtra}	% \XeLaTeX macro
 \defaultfontfeatures{Mapping=tex-text}
 \usepackage{polyglossia}	% instead of babel in xelatex
 \setdefaultlanguage{russian}
% fonts
 \setromanfont{Charis SIL}
 \setsansfont{OfficinaSansC} 
 \setmonofont{Consolas}
\fi

% several pictures in one figure
\usepackage{subfig}
% calc in TeX expressions
\usepackage{calc}
% nice pictures and plots
\usepackage{pgfplots,tikz,circuitikz}
% different libraries for pictures
\usetikzlibrary{%
  arrows,%
  calc,%
  patterns,%
  decorations.pathreplacing,%
  decorations.pathmorphing,%
  decorations.markings%
}

% colors of the hyperlinks
\hypersetup{colorlinks,%
  citecolor=blue,%
  urlcolor=blue,%
  linkcolor=red
}

\tolerance=1000
\emergencystretch=0.74cm

\newcommand{\nn}{\nonumber}
\newcommand{\pt}{\partial}
\newcommand{\eps}{\epsilon}
\newcommand{\vareps}{\varepsilon}
\newcommand{\const}{\mathrm{const}}
\newcommand{\com}[1]{{\Large{\texttt{{\color{red}(#1)}}}}}

% счётчик задач
\newcounter{notask}
\setcounter{notask}{1}

% условие без картинки
\newcommand{\task}[1]{
	\hrule
	\hbox to \textwidth {%
    	\vrule
	    \parbox[t]{0.04\textwidth}{\smallskip \centering \arabic{notask}}%
    	\vrule%
	    \hfill%
    	\parbox[t]{0.93\textwidth}{\smallskip #1 \smallskip}\hfill%
	    \vrule
	}
	\hrule
    \addtocounter{notask}{1}
    \pagebreak[2]
}

\newlength{\h}
\newsavebox{\taskbox}
\newlength{\x}
\newsavebox{\pictbox}

% условие с картинкой (картинка выравнивается по центру)
\newcommand{\taskpic}[2]{
	\savebox{\taskbox}{\parbox[t]{0.93\textwidth-4.3cm}{\smallskip #1 \smallskip}}
	\savebox{\pictbox}{\parbox[t]{4cm}{\smallskip \centering
    	\vspace{0pt} #2 \smallskip}}
	\h=\ht\taskbox
	\advance\h\dp\taskbox
	\x=\ht\pictbox
	\advance\x\dp\pictbox
	\hrule
	\hbox to \textwidth {%
		\vrule\parbox[t][\maxof{\h}{\x}][t]{0.04\textwidth}{ \smallskip
    		\centering \arabic{notask} }\vrule%
	    \hfill\parbox[t][\maxof{\h}{\x}][t]{0.93\textwidth-4.3cm}{\smallskip #1
    	    \smallskip}\hfill\vrule%
	    \hfill\parbox[t][\maxof{\h}{\x}][c]{4cm}{\hfil #2 \hfil}\hfill\vrule
	}
	\hrule
	\addtocounter{notask}{1}
	\pagebreak[2]
}

\tikzset{>=latex, interface/.style = {postaction = {draw, decorate,
	    decoration = {border, angle = 45, amplitude = 0.2cm, segment length=1.4323mm}}},%
		spring/.style={decorate,decoration={snake,amplitude=1mm,
        segment length=2mm},thick}}


\begin{document}

\thispagestyle{empty}
\parindent=5mm
\righthyphenmin=2
\begin{center}

\phantom{Бухалкин апанас}

\vfill
\LARGE{\textsc{XVII Летняя Физическая Школа}}\\
\Large{\textsc{15 июля -- 4 августа 2011}}\\[1cm]
\Large{\textit{Сборник материалов}}\\[2cm]
\includegraphics[width=6cm]{logo.pdf}
\vfill

% \small{\textsf{Санкт-Петербург}\\
% \textsf{2011}}
\end{center}

\clearpage

\section{Участники ЛФШ.}
\label{sec:participants}

\subsection{8 и 9 класс.}
\label{sec:8and9}

\begin{table}[ht]
  \begin{minipage}[t]{0.45\linewidth}\centering
    \begin{center}
      \textit{8 класс}
    \end{center}
    \begin{tabular}[h]{|c|c|c|c|}
      \hline
      \textbf{№} & \textbf{Фамилия, имя} & \textbf{Шк.} & \textbf{Д/б}\\
      \hline
      1 & Артамонов Николай & 239 & 47\\ \hline
      2 & Багиров Фарид & 533 & Д1 \\ \hline
      3 & Баринов Фёдор & 239 & Д3 \\ \hline
      4 & Богомолов Егор &  92 & ПО \\ \hline
      5 & Бомов Фёдор & 292 & 56 \\ \hline
      6 & Галкин Георгий & --- & --- \\ \hline
      7 & Гера Станислава & 534 & Д3 \\ \hline
      8 & Грибакин Борис & 74 & 50 \\ \hline
      9 & Громов Даниил & Ун. & 24 \\ \hline
      10 & Гуменюк Виталий & 92 & 66 \\ \hline
      11 & Жибарев Георгий & 533 & Д3 \\ \hline
      12 & Йона Андрей & 152 & --- \\ \hline
      13 & Кокурушников Тимофей & 239 & Д2 \\ \hline
      14 & Коробов Артём & 2 & 48 \\ \hline
      15 & Кутимский Максим & 239 & --- \\ \hline
      16 & Лебедев Александр & --- & 14 \\ \hline
      17 & Мастеров Роман & 366 & 58 \\ \hline
      18 & Морозов Дмитрий & 239 & Д2 \\ \hline
      19 & Олейник Дарья & 239 & 53 \\ \hline
      20 & Охотников Артём & 239 & Д2 \\ \hline
      21 & Петров Степан & 30 & Д2 \\ \hline
      22 & Приходько Алексей & 366 & ПО \\ \hline
      23 & Родионова Анна & 239 & 54 \\ \hline
      24 & Рутковский Илья & 78 & 46 \\ \hline
      25 & Сокольский Станислав & 366 & 33 \\ \hline
      26 & Староказников Александр & 344 & ПО \\ \hline
      27 & Трофимов Даниил & 101 & Д3 \\ \hline
      28 & Уланова Арина & 533 & Д2 \\ \hline
      29 & Усачёва Мария & 295 & 23 \\ \hline
      30 & Ходунов Павел & 239 & Д1 \\ \hline
      31 & Шубин Григорий & 550 & 42 \\ \hline
    \end{tabular}
  \end{minipage}
  \hfill
  \begin{minipage}[t]{0.45\linewidth}\centering
    \begin{center}
      \textit{9 класс}
    \end{center}
    \begin{tabular}[h]{|c|c|c|c|}
      \hline
      \textbf{№} & \textbf{Фамилия, имя} & \textbf{Шк.} & \textbf{Д/б}\\
      \hline
      1 & Агафонов Игорь & 533 & ПО  \\ \hline
      2 & Андреев Константин & ФТШ & Д2  \\ \hline
      3 & Беляков Михаил & 239 & Д1  \\ \hline
      4 & Блехштейн Максим & 239 & Д3  \\ \hline
      5 & Буренев Иван & ФТШ & Д3  \\ \hline
      6 & Вахренев Роман & 239 & 44  \\ \hline
      7 & Глезеров Евгений & 73 & ---  \\ \hline
      8 & Гордеева Людмила & ФТШ & 16  \\ \hline
      9 & Гуцол Ксения & ФТШ & ПО  \\ \hline
      10 & Затылкин Павел & 239 & 47  \\ \hline
      11 & Затылкин Кирилл & 239 & 40  \\ \hline
      12 & Иванов Владислав & 177 & Д3  \\ \hline
      13 & Киселев Егор & 239 & Д3  \\ \hline
      14 & Коско Софья & 239 & Л  \\ \hline
      15 & Кравченко Дмитрий & --- & 40  \\ \hline
      16 & Куксенок Даниил & 533 & ПО  \\ \hline
      17 & Лихачев Иван & 30 & 44  \\ \hline
      18 & Луцкий Георгий & ФТШ & ПО  \\ \hline
      19 & Малышева Александра & ФТШ & 39  \\ \hline
      20 & Муретова Мария & 239 & Д1  \\ \hline
      21 & Никоненко Михаил & 393 & ---  \\ \hline
      22 & Осипов Игорь & 571 & ---  \\ \hline
      23 & Портянкин Егор & ФТШ & 65  \\ \hline
      24 & Семенов Александр & ФТШ & 31  \\ \hline
      25 & Сычёв Станислав & ФТШ & Д1  \\ \hline
      26 & Хвещук Анастасия & 470 & ---  \\ \hline
      27 & Цейтина Елена & ФТШ & 41  \\ \hline
    \end{tabular}
  \end{minipage}
\end{table}

В \textbf{8 классе} работали: Н.В. Тараканов, М.В. Евтихиев (студент 3-го
курса ФТФ СПбГПУ), А.В. Лиознова (студентка 3-го курса ФТФ СПбГПУ),
И. Авдеев, А. Лиознов (студенты 1-го курса ФТФ СПбГПУ), Д. Максимова
(студентка 1-го курса ФТФ СПбГПУ), С. Богданов (студент 1-го курса ФФ
МГУ).

В \textbf{9 классе} работали: О.В. Шустова (студентка 6-го курса ФТФ СПбГПУ),
Ф. Затылкин (студент 2-го курса ФМФ СПбГПУ), А. Коротченков (студент
2-го курса ФТФ СПбГПУ), Ф. Петухов (студент 1-го курса ФТФ СПбГПУ). 

\clearpage

\subsection{10 и 11 класс.}
\label{sec:10and11}

\begin{table}[ht]
  \begin{minipage}[t]{0.45\linewidth}\centering
    \begin{center}
      \textit{10 класс}
    \end{center}
    \begin{tabular}[h]{|c|c|c|c|}
      \hline
      \textbf{№} & \textbf{Фамилия, имя} & \textbf{Шк.} & \textbf{Д/б}\\
      \hline
      1 & Бальков Андрей & 261 & 37\\ \hline
      2 & Вишняк Сергей & ФТШ & 41 \\ \hline
      3 & Водопьян Даниил & 239 & Р \\ \hline
      4 & Грудкин Антон & 239 & Д1 \\ \hline
      5 & Егоров Антон & ФТШ & Р \\ \hline
      6 & Иващенко Дмитрий & 239 & 32 \\ \hline
      7 & Капралов Николай & ФТШ & 46 \\ \hline
      8 & Конаныхин Роман & ФТШ & Д1 \\ \hline
      9 & Крюков Михаил & ФТШ & 31 \\ \hline
      10 & Лашкевич Злата & ФТШ & 24 \\ \hline
      11 & Люлина Анастасия & ФТШ & 27 \\ \hline
      12 & Максакова Мария & ФТШ & 39 \\ \hline
      13 & Мосягин Иван & 214 & 42 \\ \hline
      14 & Погодаев Илья & 610 & ПО \\ \hline
      15 & Рау Владислава & 239 & 46 \\ \hline
      16 & Серов Юрий & ФТШ & Р \\ \hline
      17 & Томп Дмитрий & ФТШ & Д3 \\ \hline
      18 & Ярковой Алексей & 557 & 24 \\ \hline
    \end{tabular}
  \end{minipage}
  \hfill
  \begin{minipage}[t]{0.45\linewidth}\centering
    \begin{center}
      \textit{11 класс}
    \end{center}
    \begin{tabular}[h]{|c|c|c|c|}
      \hline
      \textbf{№} & \textbf{Фамилия, имя} & \textbf{Шк.} & \textbf{Д/б}\\
      \hline
      1 & Балашов Александр & 30 & ---  \\ \hline
      2 & Борздун Наталья & 239 & ---  \\ \hline
      3 & Грачёв Дмитрий & М & ---  \\ \hline
      4 & Жаровов Дмитрий & 239 & ---  \\ \hline
      5 & Косицын Александр & 239 & ---  \\ \hline
      6 & Максимишин Дмитрий & 239 & ---  \\ \hline
      7 & Маслов Артём & ФТШ & Р  \\ \hline
      8 & Матюшин Георгий & ФТШ & Д1  \\ \hline
      9 & Михайлов Кирилл & 239 & ---  \\ \hline
      10 & Никитин Денис & 239 & Р  \\ \hline
      11 & Свирина Анна & 239 & ---  \\ \hline
      12 & Смирнов Иван & 214 & Д2  \\ \hline
      13 & Терехов Антон & 239 & Р  \\ \hline
      14 & Толстопятов Всеволод & 239 & ---  \\ \hline
      15 & Трофимов Павел & ФТШ & ---  \\ \hline
      16 & Чурилова Мария & ФТШ & ---  \\ \hline
      17 & Шалымов Роман & 239 & Д3  \\ \hline
    \end{tabular}
  \end{minipage}
\end{table}

В \textbf{10 классе} работали: И.А. Барыгин (к.ф.--м.н., преподаватель
ФТШ), С. Атамась, В. Коваленко (студенты 3-го курса ФТФ СПбГПУ).

В \textbf{11 классе} работали: И.Е. Шендерович, Д.О. Соколов
(аспиранты ПОМИ РАН), Д.С. Смирнов (студент 4-го курса ФТФ СПбГПУ). 

\section{Ежедневные занятия.}
\label{sec:daily}

\subsection{8 класс.}
\label{sec:daily8}

\subsubsection{Теория.}
\label{sec:th8}

\textit{Преподаватель: Н.В. Тараканов.}

\subsubsection{Эксперимент.}
\label{sec:exp8}

\textit{Преподаватели: А. Лиознова, Д. Максимова, И. Авдеев.}

\begin{enumerate}
\item Общие понятия о физическом эксперименте. Написание отчёта по
  эксперименту.
\item Составление плана местности в декартовых координатах.\\
  \textit{Оборудование:} линейка, два перпендикулярных друг другу
  ориентира (дорога и забор).
\item По КП, на которых написаны координаты следующей точки в
  декартовых или полярных координатах, пройти трассу.\\
  \textit{Оборудование:} компас, рулетка.
\item Построить и отградуировать на 1 минуту водяные или песочные
  часы. \\
  \textit{Оборудование:} песок, вода, скотч, одноразовые стаканчики,
  бутылка, пластилин, бумага.
\item Измерить скорость муравья.\\
  \textit{Оборудование:} нитка, линейка, секундомер.
\item Измерить ускорение свободного падения. \\
  \textit{Оборудование:} груз известной массы, линейка, секундомер.
\item Измерить трение покоя между линейкой и грузом. \\
  \textit{Оборудование:} камень, две линейки.
\item Исследовать зависимость пути тела от угла бросания. \\
  \textit{Оборудование:} резинка, камень, линейка.
\item Измерить величину центростремительного ускорения.\\
  \textit{Оборудование:} камень, нитка, секундомер.
\item От чего зависит период колебания маятника? \\
  \textit{Оборудование:} несколько разных грузов, нитки и секундомер.
\end{enumerate}

\subsection{9 класс.}
\label{sec:daily9}

\subsubsection{Теория.}
\label{sec:th9}

\textit{Преподаватель: О.В. Шустова.}

\begin{enumerate}
\item Системы координат. Векторы. Скорость. Масса. Импульс.
\item Равноускоренное движение. Графическое представление.
\item Потенциальные и непотенциальные силы. Работа. Потенциальная
  энергия.
\item Кинетическая энергия. Закон сохранения энергии.
\item Закон сохранения импульса. Абсолютно упругое/неупругое
  соударение.
\item Центр масс. Закон движения центра масс.
\item Движение по окружности. Момент силы, момент импульса.
\item Момент инерции. Расчёт для разных тел.
\item Уравнение моментов. 
\end{enumerate}

\subsubsection{Эксперимент.}
\label{sec:exp9}

\textit{Преподаватели: Ф. Затылкин, Ф. Петухов.}

\subsection{Смешанная группа 8--9 классов.}
\label{sec:daily89}

\subsubsection{Теория.}
\label{sec:th89}

\textit{Преподаватель: Н.В. Тараканов.}

\subsubsection{Эксперимент.}
\label{sec:exp89}

\textit{Преподаватели: С. Богданов, А. Лиознов.}

\subsection{10 класс.}
\label{sec:daily10}

\subsubsection{Теория.}
\label{sec:th10}

\textit{Преподаватель: С. Атамась.}

\subsubsection{Эксперимент.}
\label{sec:exp10}

\textit{Преподаватель: И.А. Барыгин.}\\
\begin{enumerate}
\item Измерить отношение масс двух грузов.\\
  \textit{Оборудование:} два груза, нитки, миллиметровка.
\item ВАХ нелинейного элемента.\\
  \textit{Оборудование:} источник постоянного тока, провода, реостат,
  мультиметр, лампочка.
\item Измерение скорости вытекания воды из крана.\\
  \textit{Оборудование:} линейка.
\item Измерение отношения длин ниток Y-образного маятника.\\
  \textit{Оборудование:} нитки, грузики.
\item Измерить жесткость каучукового шарика.\\
  \textit{Оборудование:} каучуковый шарик, вода, линейка, гуашь.
\item Измерить коэффициент трения линейки по столу.\\
  \textit{Оборудование:} две деревянные линейки.
\item Измерить зависимость мощности теплоотдачи от разницы температур.\\
  \textit{Оборудование:} горячая вода, мерный стакан или линейка,
  термометр для воды, часы, сосуд.
\end{enumerate}

\subsection{11 класс.}
\label{sec:daily11}

\subsubsection{Теория.}
\label{sec:th11}

\textit{Преподаватели: И.Е. Шендерович, Д.О. Соколов.}

\begin{enumerate}
\item Введение. Требующие объяснения экспериментальные факты.
\item Векторный анализ. Циркуляция и поток. Градиент, дивергенция и
  ротор. Теорема Стокса и теорема Гаусса--Остроградского.
\item Электростатика. Теорема Гаусса. Потенциал электростатического
  поля.
\item Переход к движущимся зарядам. Уравнение неразрывности
  электрического заряда, его физический смысл.
\item Уравнение Максвелла для ротора магнитного
  поля. Магнитостатика. Закон Био~--~Савара~--~Лапласа, закон
  Ампера. Примеры: поле кругового тока, поле соленоида, поле длинного
  провода.
\item Ток смещения, его физический смысл. Задача о разрядке
  конденсатора.
\item Электромагнитные волны. Закон индукции Фарадея как необходимое
  условие для их существования. Сила Лоренца. Примеры.
\item Физика индукции. Примеры возникновения ЭДС индукции.
\item Движущееся электромагнитное поле. Скорость волны и скорость
  света.
\item Пример полного решения уравнений Максвелла: переменное поле в
  цилиндрическом конденсаторе. Функция Бесселя.
\item Применение уравнений Максвелла к электрическим цепям. Законы
  Кирхгофа. Примеры расчёта схем.
\item Введение в теорию относительности. Инерциальные системы
  отсчета. Противоречия преобразований Галилея и уравнений Максвелла
  (на основе задачи о пластинах). Постулат об инвариантности скорости
  света. Введение системы аксиом.
\item Интервал. Физический смысл интервала. Интервал, как аналог
  расстояния. Инвариантность интервала.
\item Повороты. Преобразования Лоренца для координат. <<Парадоксы>>
  (неинерциальные системы отсчета, задача о трубе).
\item Преобразования Лоренца для скоростей. Сравнение преобразований
  Лоренца и преобразований Галилея (малые скорости,
  некоммутативность). Простые задачи на преобразование Лоренца.
\item Обзорный рассказ про преобразование для полей. Схема
  строго рассуждения (введение векторного потенциала).
\end{enumerate}

\subsubsection{Эксперимент.}
\label{sec:exp11}

\textit{Преподаватель: Д.С. Смирнов.}\\

Идеи большинства экспериментов взяты из книги С.Д. Варламова,
А.Р. Зильбермана, В.И. Зинковского <<Экспериментальные задачи на
уроках физики и физических олимпиадах>>.

\begin{enumerate}
\item Найти отношение массы монеты к массе листа миллиметровки.\\
  \textit{Оборудование:} монета, миллиметровка.
\item Найти закон распределения времени, отмеряемого внутренними
  часами. Экспериментатор должен сто раз отмерить по внутренним часам
  10 секунд и сравнить результат с истинным.
\item Определить плотность и массу куска пластилина.\\
  \textit{Оборудование:} мерный стакан, миллиметровка.
\item Измерить расстояние между бороздками CD.\\
  \textit{Оборудование:} лазерная указка известной длины волны,
  линейка.
\item Измерить и объяснить ВАХ лампочки накаливания.\\
  \textit{Оборудование:} батарейка 9В, реостат 10КОм, провода, мультиметр. 
\item Определить схему и номиналы элементов в чёрном ящике.\\
  \textit{Оборудование:} мультиметр. 
\item Определить теплоту плавления припоя.\\
  \textit{Оборудование:} свечка, мерный стакан, термопара, вода. 
\item Определить отношение частот всех мод двойного
  маятника.\\
  \textit{Оборудование:} пластилин, нитки, секундомер.
\end{enumerate}

\section{Факультативы.}
\label{sec:spec}

\begin{enumerate}
\item Квантовая механика.\\[0.2cm]
  \textit{Преподаватель: М.В. Евтихиев.}\\[0.2cm]
  История возникновения, формула Планка и фотоэффект. Атом. Модели
  атома, падение заряда на ядро и постулаты Бора. Дифракция и
  интерференция. Опыт по дифракции электронов, волна де
  Бройля. Амплитуда вероятности, принцип неопределенности.
  
\item Статистическая физика. \\[0.2cm]
  \textit{Преподаватель: И.Е. Шендерович.}\\[0.2cm]
  Задачи статистической физики. Вероятности в
  физике. Флуктуации. Связь вероятности флуктуации и минимальной
  работы в термодинамических системах. Распределение газа в закрытом
  сосуде: пример нормального распределения. Понятия среднего и
  среднеквадратичного. Статвес. Энтропия. Распределение
  Гиббса. Распределение Максвелла.
  
\item Явления переноса. \\[0.2cm]
  \textit{Преподаватель: И.А. Барыгин. }\\[0.2cm]
  Явление диффузии. Задача о пьяном матросе. Связь коэффициента
  диффузии с длиной свободного пробега. Уравнение диффузии и его
  функция Грина. Теплопроводность. Температуропроводность. Вязкость,
  динамическая и кинематическая.
\item Радуга. \\[0.2cm]
  \textit{Преподаватель: Д.С. Смирнов.}\\[0.2cm]
  Каков размер радуги? В какой последовательности в ней идут цвета?
  Где небо светлее: внутри радуги или снаружи? Сколько колец радуг
  можно увидеть одновременно? Какая последовательность цветов во всех
  этих кольцах? Поляризована ли радуга?
\item Основы астрофизики. \\[0.2cm]
  \textit{Преподаватель: В. Коваленко.}\\[0.2cm]
  Электромагнитное излучение. Понятие о фотометрии.Ослабление света
  при прохождении сквозь вещество. Понятие спектра. Тепловое
  излучение. Спектр чернотельного излучения. Формула Планка,
  приближения Рэлея--Джинса и Вина. Закон смещения Вина. Спектры
  звезд. Спектральные линии. Эффект Доплера. Смещение спектральных
  линий. Расширение спектральных линий. Спектральные приборы. Методы
  анализа наблюдений. Солнце. Общие сведения. Химический состав и
  спектр. Термоядерные реакции. Нейтринные детекторы на
  Земле. Фотосфера. Общие сведения о звездах. Спектры и
  светимости. Показатель цвета.  Спектральные классы, связь с
  температурой. Абсолютная звездная величина и светимость. Методы
  определения светимости и других характеристик далеких
  звезд. Статистические зависимости между основными
  характеристиками. Диаграмма Герцшпрунга--Рессела. Главная
  последовательность.
\item Тензоры в физике. \\
  \textit{Преподаватель: О.А. Хромов.}
\item Введение в теоретическую механику. \\[0.2cm]
  \textit{Преподаватель: Д.О. Соколов.}\\[0.2cm]
  Введение. Мотивация. Функция Лагранжа. Примеры. Метод
  вариаций. Статика. Форма верёвки в поле тяжести. Движение в поле
  тяжести. Законы сохранения: энергия, импульс, момент импульса.
\end{enumerate}

\clearpage
\section{Материалы физических боёв.}
\label{sec:battles}

\subsection{Физбой 9 класса.}
\label{sec:battle9}

\subsubsection{Полуфинал.}
\label{sec:battle9s}

\setcounter{notask}{1}

\task{Космический корабль начинает двигаться прямолинейно с
  ускорением, изменяющимся во времени так, как показано на графике
  (см. рис.). Через какое время корабль удалится от исходной точки в
  положительном направлении на максимальное расстояние? Начальная
  скорость корабля равна нулю.}

\begin{center}
  \begin{tikzpicture}
    \draw[->] (0,0) -- (0,2.5) node[right] {$a, \text{ м/c}^2$};
    \draw[->] (-0.5,1) -- (10.5,1) node[above] {$t$, с};
    \draw[dashed,blue] (0,2) -- (10.5,2);
    \draw[dashed,blue] (0,0.5) -- (10.5,0.5);
    \draw[very thick] (0,2) -- (0.5,2) -- (0.5,0.5) -- (1,0.5) --
    (1,2) -- (1.5,2) -- (1.5,0.5) -- (2.5,0.5) -- (2.5,2) -- (3,2) --
    (3,0.5) -- (4.5,0.5) -- (4.5,2) -- (5,2) -- (5,0.5) -- (7,0.5) --
    (7,2) -- (7.5,2) -- (7.5,0.5) -- (10,0.5);
    \foreach \y in {0.5,1,1.5,2} {
      \draw (-0.1,\y) -- (0.1,\y);
    }
    \foreach \x in {0.5,1,1.5,...,10} {
      \draw (\x,0.9) -- (\x,1.1);
    }
    \draw (-0.3,0.5) node {\small{-1}};
    \draw (-0.3,2) node {\small{2}};
    \draw[white,fill=white] (2.3,0.6) rectangle (2.7,0.9)
    node[midway,black] {\small{5}};
    \draw[white,fill=white] (4.8,0.6) rectangle (5.2,0.9)
    node[midway,black] {\small{10}};
    \draw[white,fill=white] (7.3,0.6) rectangle (7.7,0.9)
    node[midway,black] {\small{15}};
      \end{tikzpicture}
\end{center}

\taskpic{На гладком горизонтальном столе лежат, касаясь друг друга, две
  одинакового размера шайбы 1 и 2, радиус которых равен $R$. Шайбы
  соединены друг с другом с помощью тонкой легкой нити (см. рис., вид
  сверху). Длина нити $L = 2R$. Нить начали тянуть в горизонтальном
  направлении с постоянной силой $F$. Найдите силу, с которой шайбы
  будут давить друг на друга, когда их движение установится. Сила $F$
  приложена в середине нити. Трение можно считать малым. Рассмотрите
  два случая: 1) шайбы имеют одинаковую массу; 2) масса одной шайбы в
  два раза больше массы другой.}{
  \begin{tikzpicture}
    \draw[very thick] (0,0) circle (0.75cm);
    \draw[very thick] (0,1.5) circle (0.75cm);
    \draw[thick] (0,1.5) ++(-20:0.75) -- (1.5,0.75);
    \draw[thick] (0,0) ++(20:0.75) -- (1.5,0.75);
    \draw[very thick,->,blue] (1.5,0.75) -- (2.5,0.75) node[midway,above] {$\vec{F}$};
  \end{tikzpicture}
}

\task{Электрическая цепь составлена из семи последовательно
  соединенных резисторов: $R_1$ = 1 кОм, $R_2$ = 2 кОм, $R_3$ = 3 кОм,
  $R_4$ = 4 кОм, $R_5$ = 5 кОм, $R_6$ = 6 кОм, $R_7$ = 7 кОм и четырех
  перемычек. Входное напряжение U = 53,2 В. Укажите, в каком из
  резисторов сила тока минимальна. Найдите эту силу тока. В каком из
  резисторов сила тока максимальна? Найдите ее.}

\begin{center}
  \begin{circuitikz}
    \draw (0,0) to[generic,l_=$R_1$] (2,0) to[generic,l_=$R_2$] (4,0)
    to[generic,l_=$R_3$] (6,0) to[generic,l_=$R_4$] (8,0)
    to[generic,l_=$R_5$] (10,0) to[generic,l_=$R_6$] (12,0)
    to[generic,l_=$R_7$] (14,0);
    \draw (0,0) -- (0,-1.5) -- (6,-1.5);
    \draw (14,0) -- (14,-1.5) -- (8,-1.5);
    \draw[black,fill=white] (6,-1.5) circle (0.1cm);
    \draw[black,fill=white] (8,-1.5) circle (0.1cm);
    \draw (7,-1.5) node {$U$};
    \draw[thick,looseness=0.5] (0.2,0) to[out=80,in=100] (8,0);
    \draw[thick,looseness=0.5] (4,0) to[out=80,in=100] (12,0);
    \draw[thick,looseness=0.5] (4,0) to[out=80,in=100] (12,0);
    \draw[thick,looseness=0.5] (2,0) to[out=-80,in=-100] (10,0);
    \draw[thick,looseness=0.6] (6,0) to[out=-80,in=-100] (13.8,0);
  \end{circuitikz}
\end{center}


\clearpage
\subsubsection{Финал.}
\label{sec:battle9f}

\setcounter{notask}{1}

\task{Над обрывом установлено орудие, позволяющее вести огонь в любом
  направлении. Снаряды имеют начальную скорость $v$. На расстоянии $l$
  от орудия под углом $\varphi$ к горизонту завис воздушный шар. Известно, что
  шар находится достаточно далеко от орудия --- так, что снаряды в него
  не попадают. Обстрел стали производить снарядами, которые взрываются
  через время $T$ после выстрела. Под каким углом к горизонту следует
  стрелять, чтобы снаряды взрывались как можно ближе к шару? Ускорение
  свободного падения $g$.}

\taskpic{Массивная доска AB скользит со скоростью $u$ по гладкой
  горизонтальной поверхности. Из точки С той же поверхности
  одновременно вылетают две легкие шайбы. Первая шайба скользит по
  поверхности в направлении $CC_1$ параллельно доске АВ со скоростью
  $v_1$, вторая скользит со скоростью $v_2$ под углом $\alpha$ к
  $CC_1$. Через некоторое время шайбы сталкиваются в точке
  D. Определите скорости шайб $v_1$ и $v_2$ до столкновения, если
  известно, что время от начала движения шайб до их столкновения в $n$
  раз превышает время от начала движения шайб до столкновения второй
  шайбы с доской. При ударе шайбы о доску потерь энергии не
  происходит.}{
  \begin{tikzpicture}
    \draw[fill=gray] (0,0) rectangle (0.5,4);
    \draw (0.5,0) node[right] {$A$};
    \draw (0.5,4) node[right] {$B$};
    \draw[thick,blue,->] (0.5,2) -- (1,2) node[midway,above]
    {$\vec{u}$};
    \draw[thick] (3,0) node[right] {$C$} -- (3,4) node[right] {$C_1$};
    \draw[very thick,->,blue] (3,0.5) -- (3,1.5) node[right]
    {$\vec{v}_1$};
    \draw[very thick,->,blue] (3,0.5) -- ++(150:1.2cm) node[above]
    {$\vec{v}_2$};
    \draw[blue] (3,0.85) arc (90:150:0.35cm);
    \draw (2.8,1) node[blue] {$\alpha$};
    \draw[fill=black] (3,3) circle (0.05cm) node[right] {$D$};
  \end{tikzpicture}
}

\taskpic{Тонкостенный цилиндр катится по горизонтальной поверхности
  без проскальзывания со скоростью $v_0 = 6$ м/с. Коэффициент трения
  между цилиндром и поверхностью равен $\mu= 0,2$. Цилиндр сталкивается с
  вертикальной гладкой стенкой и упруго отражается от нее. Определите
  путь, пройденный цилиндром до остановки.}{
  \begin{tikzpicture}
    \draw[thick,interface] (3,2.5) -- (3,0) -- (0,0);
    \draw[very thick] (1,1) circle (1cm);
    \draw[thick,blue,->] (1,1) -- (2.5,1) node[above] {$\vec{v}_0$};
  \end{tikzpicture}
}

\taskpic{ По реке со скоростью $v$ плывут мелкие льдины, которые
  равномерно распределяются по поверхности воды, покрывая ее $n$-ю
  часть. В некотором месте реки образовался затор. В заторе льдины
  полностью покрывают поверхность воды, не нагромождаясь друг на
  друга.  Какая сила действует на 1 м ледяной границы между водой и
  сплошным льдом в заторе со стороны останавливающихся льдин?
  Плотность льда $\rho= 0,91 \cdot 10^3 \text{ кг/м}^3$; толщина $h =
  20$ см; скорость реки $v = 0,72$ км/ч; плывущие льдины покрывают $n =
  0,1$ часть поверхности воды. }{
  \begin{tikzpicture}
    
  \end{tikzpicture}
}

\clearpage
\subsection{Физбой 10 класса.}
\label{sec:battle10}

\setcounter{notask}{1}

\task{С какой силой отталкиваются грани равномерно заряженного по
  поверхности правильного тетраэдра? Длина ребра тетраэдра $a$, заряд
  каждой грани $q$.}

\taskpic{Участок цепи постоянного тока состоит из трех одинаковых
  вольтметров и двух одинаковых амперметров. Показания вольтметров
  V$_1$ и V$_2$ равны $U_1 = 6$~B, $U_2 = 4$~B. Как вы полагаете, что
  показывает третий вольтметр?}{
  \begin{tikzpicture}
    \draw[thick] (0,0.5) -- (0.5,0.5);
    \draw[thick] (0.5,0.5) -- (0.5,2) -- (1.45,2) (2.05,2) -- (2.5,2) --
    (2.5,1) -- (2.7,1);
    \draw[thick] (1.75,2) circle (0.3cm) node[blue] {$V_1$};
    \draw[thick] (3,1) circle (0.3cm) node[blue] {$A_2$};
    \draw[thick] (3.3,1) -- (3.5,1) -- (3.5,0.5) -- (4,0.5);
    \draw[thick] (0.5,1) -- (0.7,1);
    \draw[thick] (1,1) circle (0.3cm) node[blue] {$A_1$};
    \draw[thick] (2.6,1) -- (2.3,1);
    \draw[thick] (2,1) circle (0.3cm) node[blue] {$V_2$};
    \draw[thick] (1.3,1) -- (1.7,1);
    \draw[thick] (0.5,0.5) -- (0.5,0) -- (1.5,0) (2.1,0) -- (3.5,0) --
    (3.5,1);
    \draw[thick] (1.8,0) circle (0.3cm) node[blue] {$V_3$};
  \end{tikzpicture}
}

\task{Ракета пришельцев стартует с поверхности Земли и практически
  мгновенно набирает постоянную вертикальную скорость $v =
  800$~м/с. Злодеи обстреливают ракету из пушки, которая находится в
  $l=40$~км от стартовой площадки. Выстрел производится в момент
  старта. Могут ли злодеи попасть в ракету? Начальная скорость снаряда
  равна $u = 1000$~м/с. Сопротивлением воздуха пренебречь.}

\task{Узкий пучок протонов налетает на шар радиуса $a$. Прицельное
  расстояние (расстояние от центра шара до прямой, на которой лежит
  начальная скорость протонов) $b$, начальная энергия протонов
  $E$. Найти установившийся заряд шара. Заряд протона $e$, масса $m$.}

\taskpic{Кубик массы $M$ стоит на горизонтальной поверхности. Его
  касается кубик массы $m$, висящий на невесомой нерастяжимой
  нити. Нить составляет угол $\alpha$ с вертикалью. В начальный момент
  кубики неподвижны. Определите ускорения кубиков в начальный
  момент. Трением пренебречь. Считайте, что кубики не поворачиваются
  вокруг своей оси.}{
  \begin{tikzpicture}
    \draw[thick,interface] (4,0) -- (0,0);
    \draw[thick] (1.5,0) rectangle (3.5,2) node[midway,blue] {$M$};
    \draw[thick] (0.75,1.25) rectangle (1.5,2) node[midway,blue]
    {$m$};
    \draw (1.12,2) -- (2,4);
    \draw[fill=black] (2,4) circle (0.04cm);
    \draw[dashed] (2,4) -- (2,3);
    \draw[blue] (2,3.25) arc (270:246:0.75);
    \draw[blue] (1.8,3.03) node {$\alpha$};
  \end{tikzpicture}
}


\task{На достаточно удаленные предметы смотрят через собирающую линзу
  с фокусным расстоянием $F=9$~см, располагая глаз на расстоянии
  $a=36$~см от линзы. Оцените минимальный размер экрана, который нужно
  расположить за линзой так, чтобы он перекрыл все поле
  изображения. Считайте, что радиус зрачка равен $r=1{,}5$~мм.}


\clearpage
\subsection{Физбой 11 класса.}
\label{sec:battle11}

\setcounter{notask}{1}
\taskpic{ Два одинаковых кубика с ребром $H$ и массой $2.5m$ каждый стоят
  почти соприкасаясь гранями на гладкой горизонтальной
  поверхности. Сверху на них аккуратно кладут шар массы $m$ и радиуса
  $R$, и он начинает смещаться вертикально вниз, раздвигая кубики в
  стороны. Найти скорость шара непосредственно перед ударом о
  горизонтальную поверхность. Начальная скорость шара пренебрежимо
  мала.  }{
  \begin{tikzpicture}
    \draw[very thick,interface] (4,0) -- (0,0);
    \draw[very thick] (0.5,0) rectangle (1.5,1);
    \draw[very thick] (2.5,0) rectangle (3.5,1);
    \draw[very thick] (2,1.5) circle (0.7cm);
    \draw[blue,<->] (0.75,0) -- (0.75,1) node[midway,right] {$H$};
    \draw[blue,->] (2,1.5) node[above=0.1cm] {$R$} -- ++(40:0.7cm) ; 
  \end{tikzpicture}
}


\taskpic{ За линзой на расстоянии $\ell = 4$ см (больше фокусного)
  расположено перпендикулярно главной оптической оси плоское
  зеркало. Перед линзой, также перпендикулярно главной оптической оси,
  расположен лист клетчатой бумаги. На этом листе получают изображение
  его клеток при двух положениях листа относительно линзы. Эти
  положения отличаются на $L = 9$ см. Определить фокусное расстояние
  линзы.  }{
  \begin{tikzpicture}
    \draw[very thick] (0,0) -- (3.5,0);
    \draw[very thick] (1,1.1) -- (1,-1.1);
    \draw[very thick,<->] (2.5,1.1) -- (2.5,-1.1);
    \draw[very thick,interface] (3.5,1.1) -- (3.5,-1.1);
    \draw[blue,<->] (2.5,-1.2) -- (3.5,-1.2) node[midway,below] {$\ell$};
  \end{tikzpicture}
}

\task{ Вагон массой $M$ и длиной $L$ может без трения двигаться по
  рельсам. Он заполнен газом и разделен пополам подвижной невесомой
  вертикальной перегородкой. Вначале температура газа равна $T$. В
  правой половине включают нагреватель и доводят температуру газа до
  $2T$, в левой части температура остается прежней. Найти перемещение
  вагона, если масса всего газа равна $m$.  }

\task{ Звезда массы $M$ и радиуса $r$ образовалась из однородного
  облака газа с молярной массой $\mu$ радиуса $R$, исходно имевшего
  температуру $T$. Считая звезду также однородной, определите
  среднеквадратичное значение ее угловой скорости.  }

\task{ Радиусы кривизны двух одинаковых, слипшихся друг с другом,
  мыльных пузырей равны $R$. После того, как перегородка лопнула,
  образовался один пузырь радиусом $R_1$. Коэффициент поверхностного
  натяжения мыльного раствора $\sigma$. Найти атмосферное давление.  }

\taskpic{ Имеется равномерно заряженная диэлектрическая
  сфера. Известно, что, если ее разрезать пополам, то <<половинки>>
  будут расталкиваться с силой $F_1$. Если разрезать пополам одну из
  половинок (вдали от второй), то получившиеся <<четвертинки>> будут
  расталкиваться с силой $F_2$. И, наконец, если разрезать пополам
  одну из <<четвертинок>> (вдали от оставшихся частей сферы) на
  <<восьмушки>>, то они будут расталкиваться с силой $F_3$. Найти
  силу, с которой будут расталкиваться <<восьмушки>>, если их
  поместить так, как показано на рисунке.  }{\begin{tikzpicture}

    \def\size{50pt}


    \draw [white,fill = gray]
		(0,0) -- (\size, 0) arc (0:90:\size) -- cycle;
    \draw [white,fill = gray]
		(0,0) -- (-\size, 0) arc (180:270:\size) -- cycle;

    \fill [fill = white] (0, \size) arc (90:270:\size / 5 and \size);
    \fill [fill = white, dashed]
    	(0, -\size) arc (-90:90:\size / 5 and \size);
    \draw [fill = white, dashed,thick]
    	(\size, 0) arc (0:180:50pt and \size / 5);
    \draw [fill = white,thick] (-\size, 0) arc (180:360:50pt and \size / 5);
	\draw [thick] (0, \size) arc (90:270:\size / 5 and \size);
    \draw [dashed,thick] (0, -\size) arc
    	(-90:90:\size / 5 and \size);
    \draw[thick] (0, 0) circle (\size);
\end{tikzpicture}}




\end{document}
